Chain-of-thought (CoT) prompting is widely believed to make language models ``think alike,'' but the effect of CoT on output convergence has not been directly studied. We conduct the first controlled comparison of output convergence with and without CoT, measuring both within-model consistency (across multiple samples from the same model) and cross-model agreement (across models of different sizes). We prompt three models from the GPT-4.1 family on 90 questions spanning mathematical reasoning (GSM8K), commonsense reasoning (CommonsenseQA), and science (ARC-Challenge), collecting 2,700 responses under direct and zero-shot CoT conditions. We find that CoT's effect on convergence is \textbf{task-dependent and bidirectional}: on math tasks, CoT dramatically increases both within-model and cross-model agreement (Cohen's $d$ up to $+2.39$, $p < 0.001$), while on multiple-choice reasoning tasks, CoT \textbf{decreases} convergence (Cohen's $d$ up to $-0.91$, $p = 0.001$). The largest cross-model effects appear between models of different capability levels---CoT transforms near-random agreement ($15.5\%$) into near-perfect agreement ($94.0\%$) on math. These results show that CoT does not uniformly homogenize model outputs; rather, it acts as a convergence mechanism only when tasks admit a unique, derivable answer, with implications for ensemble design, model evaluation, and understanding AI output diversity.
